\section[Naissance de la bibliothèque]{Une bibliothèque née de l'université de droit}
L'histoire de la BIU Cujas remonte au XVe siècle où l'on mentionne pour la première fois une  \enquote{bibliothèque de l'Ecole de droit} en 1475. Cujas continua d'évoluer jusqu'à s'imposer devant une bibliothèque de référence en droit grâce à \enquote{la richesse de ses collections, l'étendue des services proposés, les missions particulières qui lui sont attribuées, ainsi que l'expérience acquise par son personnel}.En 1806, la Faculté de droit de Paris réouvre ses portes.

L'année 1829 marque l'ouverture officielle de la bibliothèque Cujas. Elle est crée comme bibliothèque de la Faculté de droit de Paris, héritière de la bibliothèque de l'Ecole de droit de l'Université de paris. Elle permet aux étudiants de préparer leurs thèses et autorise en 1831 tous les étudiants de droit de la capitale à pouvoir consulter ces collections. 

En 1876, Paul Viollet est nommé comme le premier bibliothécaire professionnel à la tête de la bibliothèque. Sous sa direction, entre 1876 et 1914, la bibliothèque Cujas connaît une forte expansion et un développement de ses collections, une augmentation de son budget et l'accroissement de son effectif.   

Depuis la signature du traité de paix, chacun pense que \enquote{l'on va pouvoir beaucoup changer l'état des choses}, mais en réalité les années 1919-1939 sont \enquote{une période noire pour les bibliothèques d'études et de recherche}. Le contexte n'est pas favorable au développement des bibliothèques, ni à une refonte de nos institutions publiques à cause de l'effondrement de la Troisième République. A la sortie de la Seconde Guerre Mondiale, la Société des Nations a conscience de l'importance que revêt la culture, crée dans une \enquote{Commission internationale de coopération intellectuelle, présidée par Henri Bergson}. Toutefois, l'appui financier venant à manquer et la multiplication des bibliothèques et des étudiants rendent difficile la mise en place des réformes, ainsi que son organisation. \enquote{Depuis le décret du 9 novembre 1946, il n'existe plus une bibliothèque de l'université de Paris mais des bibliothèques facultaires, auxquelles sont théoriquement rattachées les bibliothèques de laboratoires et dinstitutions de la faculté correspondante}.

Après les événements de mai 1968, la loi du 12 novembre 1968 est promulguée par le ministre Edgar Faure, dont la volonté est de réformer l'université accordant une autonomie de gestion et d'organisation. La conséquence indirecte de la loi Faure est qu'elle menace de disperser les livres précieux.  En 1970, la partition de l'Université entraîne la disparition de l'ancienne faculté, affectant les autres bibliothèques membres comme la bibliothèque Cujas. En effet, le gouvernement et l'État visent à éviter la dispersion des fonds particuliers. Ils possèdent une importance majeure pour la recherche, dans la mesure où leur \enquote{rôle national, rendait difficile leur attribution à une université plutôt qu'à une autre}. La décision a été prise de constituer des bibliothèques interuniversitaires. La bibliothèque Cujas se trouve \enquote{intégrée dans la bibliothèque interuniversitaire A}, changeant ainsi de statut pour devenir la Bibliothèque universitaire Cujas. 

En 1979, la Bibliothèque interuniversitaire Cujas est rattaché administrativement à l'Université Paris I, Panthéon-Sorbonne. Puis, grâce à une convention signée cette même année, elle se trouve co-gérée et conventionnée entre les universités Panthéon-Sorbonne (Paris I) et Panthéon-Assas(Paris II).
\bigskip

La BIU Cujas est spécialisée en droit, économmie et sciences politiques devenue aujourd'hui une \enquote{bibliothèque de référence nationale et internationale dans ces domaines}.
L'une des missions principales de Cujas est \enquote{d'offrir à la communauté universitaire et scientifique l'exhaustivité des ressources documentaires françaises ainsi qu'une vaste sélection de la production étrangère en sciences juridiques, économiques et politiques}. Son prestige découle de sa richesse documentaire \enquote{d'environ un million de documents (monographies, encyclopédies, ouvrages à feuillets mobiles, périodiques, thèses, mémoires) dont de nombreux ouvrages précieux}. Entre ses murs, la bibliothèque Cujas conserve \enquote{sept incunables, le plus ancien datant de 1472}.

La saturation de ces magasins est l'un des défis auxquels fait face la bibliothèque Cujas. Celle-ci est d'ailleurs loin d'être la seule face à ce phénomène, qui touche régulièrement la majorité des institutions patrimoniales. En effet, alors que les collections s'enrichissent chaque année, l'espace de stockage physique disponible, quant à lui, diminue. Cette réduction d'espace amène ainsi la bibliothèque à repenser ses modes de conservation et à concevoir d'autres moyens de diffusion pour ses collections.